\section{Egy szabadsági fokú, csillapítatlan rezgőrendszer mozgástörvénye}

\subsection*{Karakterisztikus egyenlet levezetése a próbafüggvény módszerrel}
Az egy szabadsági fokú csillapítatlan rezgőrendszer referencia mozgásegyenlete a következő alakú:
\begin{equation}
    m\ddot{x} + kx = 0 \quad \text{vagyis} \quad \ddot{x} + \omega_{n}^2 x = 0
\end{equation}
Ahol $\omega_n = \sqrt{k/m}$ (referencia saját-körfrekvencia).
A megoldást a próbafüggvény exponenciális alakjában keressük:
\begin{equation}
    x(t) = C e^{\lambda t}
\end{equation}
Ennek deriváltjai: $\dot{x}(t) = \lambda C e^{\lambda t}$, $\ddot{x}(t) = \lambda^2 C e^{\lambda t}$. Visszahelyettesítve a mozgásegyenletbe:
\begin{equation}
    \left( \lambda^2 + \omega_n^2 \right) C e^{\lambda t} = 0
\end{equation}
Mivel a nemtriviális megoldás érdekében $C e^{\lambda t} \neq 0$, a zárójelben lévő kifejezésnek kell nullának lennie. Így kapjuk meg a karakterisztikus egyenletet:
\begin{equation}
    \lambda^2 + \omega_n^2 = 0
\end{equation}

\subsection*{Karakterisztikus gyökök és komplex számsíkon való elhelyezkedésük}
A karakterisztikus egyenletből a gyökök kifejezve:
\begin{equation}
    \lambda_{1,2} = \pm i \omega_n
\end{equation}
amelyek egy tiszta képzetes konjugált gyökpárt alkotnak.
\begin{figure}[H]
    \centering
    \begin{tikzpicture}
        % axes
        \draw[->, thick] (-2, 0) -- (2, 0) node[right] {Re};
        \draw[->, thick] (0, -2) -- (0, 2) node[above] {Im};
        % roots
        \filldraw[red] (0, 1.5) circle (2pt) node[right] {$+i\omega_n$};
        \filldraw[red] (0, -1.5) circle (2pt) node[right] {$-i\omega_n$};
    \end{tikzpicture}
    \caption{A karakterisztikus gyökök elhelyezkedése a komplex számsíkon.}
\end{figure}

\subsection*{Az általános megoldás felírásának lépései}
Mivel a gyökök tiszta képzetesek ($\lambda_{1,2} = \pm i \omega_n$), az általános megoldás felírható az Euler-formulák felhasználásával a próbafüggvények lineáris kombinációjaként:
\begin{equation}
    x(t) = C_1 e^{i\omega_n t} + C_2 e^{-i\omega_n t}
\end{equation}
A fizikai értelmezhetőség (tisztán valós kitérés) miatt az Euler formula segítségével a megoldást trigonometrikus alakba konvertáljuk:
\begin{equation}
    x(t) = A \cos(\omega_n t) + B \sin(\omega_n t)
\end{equation}
vagy egyenértékű amplitúdó-fázisszög alakban definiáljuk:
\begin{equation}
    x(t) = X_0 \sin(\omega_n t + \phi) \quad \text{vagy} \quad x(t) = X_0 \cos(\omega_n t - \phi')
\end{equation}
Az eljárás utolsó lépése az $A, B$ (vagy $X_0, \phi$) konstansok értékének meghatározása a perem- vagy kezdeti feltételekből ($x(0), \dot{x}(0)$).
